\documentclass{article}

\usepackage[russian]{babel}
\usepackage[utf8]{inputenc}
\usepackage{amsthm,amsmath, amsfonts }

\usepackage[dvipdfm]{graphicx}
%если будет конвертация в pdf командой dvipdfm (но в dvi-файле картинка может быть не видна), 
%и \usepackage{graphicx} если рисунок нужно вставить в dvi (bmp может и не вставиться)

\newtheorem{defi}{Определение.}
\newtheorem{theorem}{Теорема.}

\theoremstyle{remark}
\newtheorem{remark}{Замечание.}

\begin{document}

\begin{titlepage}
\fontsize{200pt}{150pt}
\centerline
{\bf   Gr\"obner}
\vspace{1 cm}
\centerline
{ \bf Basis }
%\begin{figure}[h]
\hspace{1cm}
\includegraphics*[width=120mm, height=60mm]{1.bmp} 
%\end{figure}
\vspace{1cm}
\hspace{2cm}
{\bf  Петрова А.А.}
\vspace{1cm}
\fontsize{15pt}{10pt}
\begin{center}
{\it Москва \\ 2024 }
\end{center}
\end{titlepage}


\section
{
  Введение.
}
Формулировка задачи....

Обсуждение....
\begin{remark}\label{n1}
Этот текст выплнен........

$(x_1>x_2>=x_3)$

\end{remark}

\[
\left\{ U_i,\  i=1,\dots,k
\right\}=\prod_{1\le j\le n}\alpha_j| \
{\alpha_j\notin\bold Z}
\]

\section
{
  Предварительные факты.
}
Введем ряд обозначений
$$
A= \left(
\begin{array}{cc}
1&{\fontsize{10pt}{5pt}\bf VII}\\
5&655  \end{array}
\right)
$$
.......


\section{Основные результаты.}
\begin{theorem}\label{Main}
Пусть.....Тогда......
\begin{align}\label{eq1}
&A=1;\\ \notag
&B=2;\\ 
&{\mathcal A}_5={\mathcal Z}\notag
\end{align}
\end{theorem}


\subsection{ Леммы и доказательство теоремы.}
 Рассмотрим ряд лемм....
и докажем  теорему \ref{Main}.
\begin{proof} Рассмотрим... , 
принимая во внимание  \cite{Jo1} и  замечание
\ref{n1},
можно утверждать.....

\begin{enumerate}
\item   По лемме 1....
\item   По лемме 2 .....
\end{enumerate}
Следовательно,  .....и верно уравнение (\ref{eq1})
\end{proof}

\subsection{Алгоритм.}

\begin{defi}
Последовательность $\{a_n\}$ называется фундаментальной,
если выполняется следующее:
 $\forall\varepsilon>0\ \exists n_0=n_0(\varepsilon)$ 
такой, что
\(\forall m,n\in {\mathbb N}, m,n>n_0\) имеем \(|a_n-a_m|<\varepsilon\). 
\end{defi}

\begin{verbatim}
{v_i    $
\end{verbatim}

\begin{thebibliography}{}
\bibitem{Jac} 
Jacobi C.G.J. 
{\it De
    investigando ordine systematis aequationum differentialum
    vulgarium cujuscunque}.
Borchardt Journal f{\"u}r die reine und angewandte Mathematik,
 Ed. C.~W.~Borchardt. Berlin,
    Reimer.
N 4, LXIV (1865), p. 297-320.
Also republ.:
Jacobi C.G.J. 
 {\it Gesammelte Werke}.
V. 5      (1890),
  p. 191-216.
 Ed. Georg Reimer,
Berlin.
\bibitem{Jo1} Johnson Joseph.  {\it K\"ahler differentials and differential algebra} 
Ann. of Math., 89 (1969), p. 92-98.
\bibitem{Kolchin} Kolchin E.R. Differential algebra and algebraic groups. Academic Press. New York-London. 1973.
\bibitem{Lando} Lando B. {Jacobi's bound for the order of systems of first order   differential equations. }
Trans. Amer. Math. Soc., N 1, v. 152 (1970), p. 119-135. 
\bibitem{Rit1} Ritt J.F. {\it Jacobi's problem on the order of systems of differential equations.} Ann. of Math., 36 (1935). p.303-312.
\bibitem{Rit2} Ritt J.F.  {\it Differential Algebra,} volume XXXIII of
Colloquium
Publications. New York, American Mathematical Society, 1950.
\bibitem{To} Tomasovic T.S. {\it A Generalized Jacobi Conjecture
for Arbitrary Systems of Algebraic Differential Equations,} Th. D. dissertation, 
Columbia University, 1976.
\end{thebibliography}
\end{document}